\documentclass[a4paper,11pt]{article}
\usepackage[utf8]{inputenc}
\usepackage[T2A]{fontenc}
\usepackage[english,russian]{babel}
\usepackage{amsmath,amssymb}
\usepackage{geometry}
\geometry{margin=2.5cm}
\usepackage{booktabs}
\usepackage{listings}
\usepackage{xcolor}
\usepackage{hyperref}

\lstset{
  basicstyle=\ttfamily\small,
  keywordstyle=\color{blue},
  commentstyle=\color{gray},
  stringstyle=\color{red},
  breaklines=true,
  frame=single,
  language=C++
}

\title{Отчёт об исправлении 15-кратного расхождения \\ MBS-raw vs MBS-old}
\author{MBS-fast project}
\date{24 марта 2026}

\begin{document}
\maketitle

\section{Описание проблемы}

При сравнении нового оптимизированного кода (MBS-raw/MBS-fast) со старым бинарником
(MBS-old) в режиме фиксированной ориентации \texttt{--fixed 45 30} обнаружено
15-кратное расхождение в значениях $M_{11}$:

\begin{center}
\begin{tabular}{lrl}
\toprule
Код & $M_{11}(0^\circ)$ & Параметры \\
\midrule
MBS-old & $5.76 \times 10^5$ & \texttt{--fixed 45 30 --grid 0 25 360 25} \\
MBS-raw (новый) & $8.93 \times 10^6$ & те же + \texttt{--legacy\_sign} \\
\bottomrule
\end{tabular}
\end{center}

Отношение $\approx 15.5\times$. Флаг \texttt{--legacy\_sign} (знак Френеля) влияния
не оказал.

\section{Методика диагностики}

\subsection{Шаг 1: Инкогерентный режим}

Запуск обоих бинарников с флагом \texttt{--incoh} (мюллеровские матрицы суммируются
поштучно по лучам, без когерентного сложения амплитуд):

\begin{center}
\begin{tabular}{lcl}
\toprule
Тест & Макс. отн. разница & Вывод \\
\midrule
\texttt{--fixed 0 0 --incoh} & $< 10^{-6}$ & Поштучная дифракция идентична \\
\texttt{--fixed 0 0} (когерент.) & $\sim 5\%$ & Фазы отличаются \\
\bottomrule
\end{tabular}
\end{center}

\textbf{Вывод:} амплитуды $|J_\text{beam}|$ идентичны, расхождение только в фазах
при когерентном суммировании.

\subsection{Шаг 2: Обнаружение отсутствия поворота частицы}

Сравнение кода \texttt{TracerPO::TraceFixed} в MBS-old и MBS-raw:

\begin{lstlisting}
// MBS-old: НЕ поворачивает частицу!
void TracerPO::TraceFixed(double beta, double gamma) {
    double b = DegToRad(beta);
    double g = DegToRad(gamma);
    m_scattering->ScatterLight(b, g, outBeams);  // углы игнорируются
    m_handler->HandleBeams(outBeams, 1);
}

// MBS-raw (новый): ПОВОРАЧИВАЕТ правильно
void TracerPO::TraceFixed(double beta, double gamma) {
    double b = DegToRad(beta);
    double g = DegToRad(gamma);
    m_particle->Rotate(b, g, 0);                 // <-- ДОБАВЛЕНО
    m_scattering->ScatterLight(b, g, outBeams);
    m_handler->HandleBeams(outBeams, 1);
}
\end{lstlisting}

В \texttt{ScatteringConvex::ScatterLight} вызов \texttt{m\_particle->Rotate()}
\textbf{закомментирован}. Старый код всегда считает рассеяние для ориентации по
умолчанию, игнорируя аргументы \texttt{--fixed}.

\textbf{Это баг старого кода.} Для ориентационного усреднения
(\texttt{TracerPOTotal}) баг не проявляется, т.к. \texttt{Rotate} вызывается
перед \texttt{ScatterLight(0, 0, ...)}.

\subsection{Шаг 3: Остаточное расхождение 5\% при --fixed 0 0}

После устранения фактора поворота (тест при $\beta=0, \gamma=0$) осталось
$\sim 5\%$ расхождение в когерентном режиме. Инкогерентный --- точное совпадение.

Причина: вызов \texttt{Rotate(0, 0, 0)} в новом коде запускает
\texttt{SetDParams()}, которая инициализирует \texttt{d\_param} (4-й компонент
\texttt{Point3f::coordinates[3]}).

\section{Найденные баги}

\subsection{Баг 1: Неинициализированный d\_param}

\texttt{Point3f} имеет пустой конструктор по умолчанию:
\begin{lstlisting}
struct Point3f {
    float coordinates[4];
    Point3f() {}                    // НЕ обнуляет coordinates[3]!
    Point3f(float x, float y, float z) {
        coordinates[0] = x;
        coordinates[1] = y;
        coordinates[2] = z;         // coordinates[3] НЕ задан!
    }
};
#define d_param coordinates[3]
\end{lstlisting}

Во всех конструкторах частиц:
\begin{lstlisting}
SetDefaultNormals();   // вычисляет нормали (cx,cy,cz), НО НЕ d_param
Reset();               // копирует defaultFacets -> facets
SetDefaultCenters();   // вычисляет центры
// SetDParams() НЕ ВЫЗЫВАЕТСЯ -> d_param мусорный!
\end{lstlisting}

Только \texttt{Rotate()} вызывает \texttt{SetDParams()}.

\subsubsection*{Исправление}

Добавлен \texttt{SetDParams()} после \texttt{Reset()} во все 8 конструкторов
частиц:
\begin{itemize}
\item \texttt{Hexagonal.cpp}
\item \texttt{ConcaveHexagonal.cpp}
\item \texttt{Bullet.cpp}
\item \texttt{BulletRosette.cpp}
\item \texttt{Droxtal.cpp}
\item \texttt{HexagonalAggregate.cpp}
\item \texttt{CertainAggregate.cpp}
\item \texttt{Particle.cpp} (загрузка из файла)
\end{itemize}

\subsection{Баг 2: Смешение fast\_exp\_im и exp\_im}

В fallback-функциях дифракции (\texttt{DiffractIncline},
\texttt{DiffractInclineFast}) в одном выражении смешивались
\texttt{fast\_exp\_im} (полиномиальный sincos) и \texttt{exp\_im} (glibc sincos):

\begin{lstlisting}
// БЫЛО: смешение точностей
tmp = (fast_exp_im(kCi*p2.x) - exp_im(kCi*p1.x)) / Ci;
s += fast_exp_im(m_waveIndex*B*bi) * tmp;

// СТАЛО: единообразная точность
tmp = (exp_im(kCi*p2.x) - exp_im(kCi*p1.x)) / Ci;
s += exp_im(m_waveIndex*B*bi) * tmp;
\end{lstlisting}

Аналогично исправлены \texttt{ComputeFnJones}, \texttt{ApplyDiffraction},
\texttt{ApplyDiffractionFast}, \texttt{ApplyDiffractionFast2}.

\textbf{Примечание:} инлайн-путь HandleBeams по-прежнему использует
\texttt{fast\_sincos} повсеместно (нет смешения), что корректно и быстро.

\section{Результаты валидации}

\begin{center}
\begin{tabular}{llrl}
\toprule
Тест & Режим & Макс. отн. разн. & Примечание \\
\midrule
\texttt{--fixed 0 0} & incoh & $3.3 \times 10^{-5}$ & d\_param фикс \\
\texttt{--fixed 0 0} & coh, без Rotate & $< 10^{-6}$ & точн. совп. \\
\texttt{--fixed 0 0} & coh, с Rotate & $\sim 5\%$ & d\_param: мусор vs правильный \\
\texttt{--fixed 45 30} & old binary & \multicolumn{2}{l}{всегда = ориентация 0 (баг)} \\
\texttt{--fixed 45 30} & new binary & \multicolumn{2}{l}{корректный поворот} \\
\bottomrule
\end{tabular}
\end{center}

\section{Влияние на рабочие режимы}

\begin{center}
\begin{tabular}{lcc}
\toprule
Режим & Затронут? & Причина \\
\midrule
\texttt{--random} (TracerPOTotal) & Нет & Rotate() всегда вызывается \\
\texttt{--sobol} & Нет & То же \\
\texttt{--adaptive / --auto} & Нет & То же \\
\texttt{--fixed} (TracerPO) & \textbf{Да} & Добавлен Rotate + d\_param \\
\texttt{--montecarlo} (TracerPOTotal) & Нет & Rotate() всегда вызывается \\
\bottomrule
\end{tabular}
\end{center}

Для всех рабочих режимов ориентационного усреднения (random, sobol, adaptive, auto)
исправление не меняет результатов, т.к. \texttt{TracerPOTotal} всегда вызывает
\texttt{Rotate()} перед \texttt{ScatterLight()}, что инициализирует \texttt{d\_param}.

\section{Заключение}

\begin{enumerate}
\item Исходное 15-кратное расхождение --- не баг HandleBeams, а \textbf{сравнение
  разных ориентаций}: старый код не поворачивает частицу в режиме \texttt{--fixed}.
\item Остаточное 5\% расхождение в когерентном режиме --- от неинициализированного
  \texttt{d\_param}, исправлено во всех конструкторах.
\item Инлайн-оптимизации HandleBeams (batched sincos, precomputed edges,
  theta-coefficients) \textbf{верифицированы}: дают идентичный результат поштучно.
\item Fallback-функции приведены к единообразному использованию \texttt{exp\_im}.
\end{enumerate}

\end{document}
